% !TeX program = xelatex
% 中文 Beamer 模板，参考用户提供的经典蓝色样式。
% XeLaTeX 编译两次，或使用 Tectonic。图表仅为合成示例。
\documentclass[aspectratio=169,11pt]{beamer}
\usepackage[UTF8,fontset=fandol]{ctex}
\usepackage{amsmath,amssymb,booktabs}
\usepackage{pgfplots}
\pgfplotsset{compat=1.18}

% ===== 经典 Beamer 蓝：与参考截图主色一致 =====
\usetheme{Madrid}
\usefonttheme{professionalfonts}
\definecolor{AcademicBlue}{RGB}{51,51,178}
\definecolor{BlueDark}{RGB}{26,26,89}
\definecolor{BlueMid}{RGB}{38,38,134}
\definecolor{NoteGray}{RGB}{105,105,105}
\definecolor{GridGray}{RGB}{220,220,230}
\setbeamercolor{structure}{fg=AcademicBlue}
\setbeamercolor{palette primary}{fg=white,bg=AcademicBlue}
\setbeamercolor{palette secondary}{fg=white,bg=BlueMid}
\setbeamercolor{palette tertiary}{fg=white,bg=BlueDark}
\setbeamercolor{frametitle}{fg=white,bg=AcademicBlue}
\setbeamercolor{title}{fg=white,bg=AcademicBlue}
\setbeamercolor{subtitle}{fg=white}
\setbeamercolor{normal text}{fg=black,bg=white}
\setbeamercolor{block title}{fg=white,bg=AcademicBlue}
\setbeamercolor{block body}{fg=black,bg=AcademicBlue!7!white}
\setbeamercolor{author in head/foot}{fg=white,bg=BlueDark}
\setbeamercolor{title in head/foot}{fg=white,bg=BlueMid}
\setbeamercolor{date in head/foot}{fg=white,bg=AcademicBlue}

% 中文标题不额外加粗，保持参考图的轻重关系。
\setbeamerfont{title}{size=\Large,series=\mdseries}
\setbeamerfont{subtitle}{size=\normalsize,series=\mdseries}
\setbeamerfont{frametitle}{size=\Large,series=\mdseries}
\setbeamerfont{author}{size=\normalsize}
\setbeamerfont{institute}{size=\small}
\setbeamerfont{date}{size=\normalsize}
\setbeamersize{text margin left=4mm,text margin right=4mm}
\setbeamertemplate{navigation symbols}{}
\setbeamertemplate{headline}{}
\setbeamertemplate{enumerate items}[ball]
\setbeamertemplate{sections/subsections in toc}[ball]
\setbeamertemplate{blocks}[rounded][shadow=true]

% 参考图的三段深浅蓝页脚，封面也保留。
\setbeamertemplate{footline}{%
  \leavevmode\hbox{%
  \begin{beamercolorbox}[wd=.333333\paperwidth,ht=2.1ex,dp=.85ex,center]{author in head/foot}%
    \fontsize{5}{6}\selectfont\bfseries\insertshortauthor\quad（\insertshortinstitute）
  \end{beamercolorbox}%
  \begin{beamercolorbox}[wd=.333333\paperwidth,ht=2.1ex,dp=.85ex,center]{title in head/foot}%
    \fontsize{5}{6}\selectfont\bfseries\insertshorttitle
  \end{beamercolorbox}%
  \begin{beamercolorbox}[wd=.333334\paperwidth,ht=2.1ex,dp=.85ex]{date in head/foot}%
    \fontsize{5}{6}\selectfont\bfseries\hspace{5mm}\insertshortdate
    \hfill\insertframenumber\,/\,\inserttotalframenumber\hspace{1.5mm}
  \end{beamercolorbox}}}

\newcommand{\smallnote}[1]{\par{\scriptsize\color{NoteGray}#1\par}}
\newcommand{\bluelabel}[1]{{\color{AcademicBlue}#1}\par}
\newcommand{\spaceditem}{\setlength{\itemsep}{.65em}}

% ===== 报告信息：替换这里即可复用 =====
\title[学术汇报模板]{空间转录组学中的计算与人工智能\\学术汇报模板示例}
\subtitle{研究背景、问题建模与方法分析}
\author[汇报人]{汇报人：姓名}
\institute[单位]{某某大学\quad 某某学院 / 课题组}
\date[2026年10月21日]{2026年10月21日}
\hypersetup{pdftitle={中文学术汇报模板：经典蓝色 Beamer},pdfauthor={模板示例}}

\begin{document}

% 01 / 圆角蓝色标题框，居中姓名、单位与日期
\begin{frame}
  \titlepage
\end{frame}

% 02 / 与参考图一致的圆球编号目录
\begin{frame}{汇报提纲}
  \begin{enumerate}
    \setlength{\itemsep}{6mm}
    \color{AcademicBlue}
    \item 研究背景与动机
    \item 问题定义
    \item 研究方法
    \item 实验结果与分析
    \item 总结与展望
  \end{enumerate}
\end{frame}

\section{研究背景与动机}
% 03 / 以双栏呈现研究背景
\begin{frame}{研究背景：表达信息与空间信息}
  \begin{columns}[T,onlytextwidth]
    \begin{column}{.46\textwidth}
      \bluelabel{表达矩阵}
      \[
        X\in\mathbb{R}^{n\times p}
      \]
      \begin{itemize}\spaceditem
        \item $n$ 个观测位置
        \item $p$ 个测量特征
        \item 每行对应一个位置
      \end{itemize}
    \end{column}
    \begin{column}{.46\textwidth}
      \bluelabel{空间坐标}
      \[
        S\in\mathbb{R}^{n\times 2}
      \]
      \begin{itemize}\spaceditem
        \item 二维平面中的观测位置
        \item 邻接矩阵 $A\in\mathbb{R}^{n\times n}$
        \item $A$ 描述空间邻接关系
      \end{itemize}
    \end{column}
  \end{columns}
  \vspace{3mm}
  \begin{block}{研究动机}
    在低维表示中同时考虑表达信息与空间邻接关系。
  \end{block}
  \vspace{2mm}
  \smallnote{教学示例。变量与符号用于展示版式，尚未转换讲义内容。}
\end{frame}

\section{问题定义}
% 04 / 传统 Beamer 信息层级
\begin{frame}{问题定义：输入、输出与建模假设}
  \begin{block}{输入与输出}
    给定表达矩阵 $X$ 和空间邻接矩阵 $A$，学习低维表示
    $Z\in\mathbb{R}^{n\times d}$，其中 $d\ll\min(n,p)$。
  \end{block}
  \vspace{4mm}
  \bluelabel{建模假设}
  \begin{itemize}\spaceditem
    \item $X$ 的主要变化可以用低维表示近似。
    \item 对选定的空间图，模型鼓励相邻位置具有相近的表示。
  \end{itemize}
  \vspace{4mm}
  \bluelabel{需要检查的边界}
  \begin{itemize}
    \item 真实组织边界可能连接不同状态，需检验平滑假设是否合适。
  \end{itemize}
  \vfill
  \smallnote{以上为示例模型的设定与假设，不是实验结论。}
\end{frame}

\section{研究方法}
% 05 / 公式页
\begin{frame}{研究方法：带空间正则项的低秩表示}
  \[
    \min_{\substack{Z,W\\W^{\mathsf T}W=I_d}}
    \underbrace{\lVert X-ZW^{\mathsf T}\rVert_F^2}_{\text{表达重构}}
    +\lambda\underbrace{\operatorname{tr}(Z^{\mathsf T}LZ)}_{\text{空间正则}}
  \]
  \vspace{5mm}
  \begin{columns}[T,onlytextwidth]
    \begin{column}{.46\textwidth}
      \bluelabel{变量与参数}
      \begin{itemize}\spaceditem
        \item $Z\in\mathbb{R}^{n\times d}$：低维表示
        \item $W\in\mathbb{R}^{p\times d}$：特征载荷
        \item $\lambda\geq0$：正则项权重
      \end{itemize}
    \end{column}
    \begin{column}{.46\textwidth}
      \bluelabel{图拉普拉斯矩阵}
      \[
        L=D-A,\qquad D_{ii}=\sum_j A_{ij}
      \]
      \small 对称非负邻接矩阵下，该正则项惩罚相邻位置的表示差异。
    \end{column}
  \end{columns}
  \vfill
  \smallnote{教学模型示例，用于展示数学公式与符号说明的排版。}
\end{frame}

\section{实验结果与分析}
% 06 / 可编辑曲线图
\begin{frame}{实验结果：目标函数变化示例}
  \begin{columns}[c,onlytextwidth]
    \begin{column}{.66\textwidth}
      \begin{tikzpicture}
      \begin{axis}[
        width=.96\linewidth,height=50mm,
        xmin=0,xmax=50,ymin=0,ymax=1.05,
        xlabel={迭代次数},ylabel={归一化目标函数值},
        label style={font=\small},tick label style={font=\scriptsize},
        xtick={0,10,20,30,40,50},ytick={0,.25,.5,.75,1},
        axis lines=left,axis line style={black!60},
        tick style={black!60},grid=major,grid style={GridGray,dashed},
        legend style={at={(.97,.97)},anchor=north east,
          draw=none,fill=white,font=\scriptsize}]
        \addplot[AcademicBlue,thick,mark=*,mark size=1.6pt]
          coordinates {(0,1)(10,.64)(20,.43)(30,.32)(40,.27)(50,.25)};
        \addlegendentry{示例设置 A}
        \addplot[black!55,thick,dashed,mark=square*,mark size=1.5pt]
          coordinates {(0,1)(10,.51)(20,.32)(30,.24)(40,.21)(50,.20)};
        \addlegendentry{示例设置 B}
      \end{axis}
      \end{tikzpicture}
    \end{column}
    \begin{column}{.29\textwidth}
      \bluelabel{结果观察}
      \vspace{2mm}
      两条示意曲线随迭代下降，后期变化趋缓。
      \par\vspace{5mm}
      \bluelabel{结果解释}
      \vspace{2mm}
      \small 真实实验需说明重复次数与不确定性。
    \end{column}
  \end{columns}
  \vspace{4mm}
  \smallnote{合成数据，仅展示图表版式，不代表任何真实算法的性能。}
\end{frame}

% 07 / 三线表
\begin{frame}{实验结果：表示维度与计算开销}
  \begin{center}
    \renewcommand{\arraystretch}{1.55}
    \setlength{\tabcolsep}{8mm}
    \begin{tabular}{@{}lrrr@{}}
      \toprule
      设置 & 维度 $d$ & 重构误差 & 耗时 / s\\
      \midrule
      设置 A & 8  & 0.38 & 12.4\\
      设置 B & 16 & 0.29 & 18.7\\
      设置 C & 32 & 0.24 & 31.6\\
      \bottomrule
    \end{tabular}
  \end{center}
  \vspace{4mm}
  \begin{block}{示例解读}
    在这组合成数值中，维度越高，误差越低，耗时越长。
  \end{block}
  \vspace{3mm}
  \smallnote{合成数据。误差为无量纲示意值，耗时不来自实际测量。}
\end{frame}

\section{总结与展望}
% 08 / 总结与讨论
\begin{frame}{总结与展望}
  \begin{enumerate}\spaceditem
    \item {\color{AcademicBlue}模型表达}\\[1mm]
      同一目标函数组合表达重构项与空间正则项。
    \item {\color{AcademicBlue}参数权衡}\\[1mm]
      $\lambda$ 控制两项的相对权重，$d$ 决定表示维度。
    \item {\color{AcademicBlue}后续验证}\\[1mm]
      比较空间图的构建方式，并设计独立验证策略。
  \end{enumerate}
  \vspace{4mm}
  \begin{block}{讨论问题}
    当相邻位置跨越组织边界时，怎样调整空间约束？
  \end{block}
  \vfill
  \smallnote{教学模型的讨论提纲。真实结论需由数据与实验支持。}
\end{frame}

% 09 / 结束页，沿用封面的标题框与三段页脚
\begin{frame}
  \vfill
  \begin{beamercolorbox}[sep=8mm,center,rounded=true,shadow=true]{title}
    {\Large 谢谢，欢迎批评指正}
  \end{beamercolorbox}
  \vspace{8mm}
  \begin{center}
    汇报人：姓名\\[3mm]
    {\small 某某大学\quad 某某学院 / 课题组}
  \end{center}
  \vfill
\end{frame}

\end{document}
